\documentclass{article}
\usepackage[utf8]{inputenc}
\title{Rapporto di Ricerca: L'Impatto dei Social Media sull'Immagine Corporea negli Adolescenti}
\author{}
\date{}

\begin{document}

\maketitle

\section{Domanda di Ricerca}
In che modo l'esposizione a contenuti sui social network può essere associata alla percezione dell'immagine corporea negli adolescenti?

\section{Metodo di Consultazione}
La ricerca è stata condotta utilizzando OpenAlex per identificare documenti pertinenti. Sono state utilizzate le seguenti parole chiave: "body image", "social media", "adolescents". Sono stati selezionati i documenti più rilevanti per esplorare i temi e le sotto-domande del piano di ricerca.

\section{Risultati}
Sono stati consultati i seguenti documenti:

1. **Social Media Effects Regarding Eating Disorders and Body Image in Young Adolescents**  
   \textit{ID:} \href{https://openalex.org/W4394983186}{W4394983186}  
   Rivede come i social media influenzano l'immagine corporea e i disturbi alimentari negli adolescenti, evidenziando che l'esposizione a modelli di bellezza irraggiungibili promuove insoddisfazione corporea (Suhag \& Rauniyar, 2024).

2. **Processing Body Image on Social Media: Gender Differences in Adolescent Boys’ and Girls’ Agency and Active Coping**  
   \textit{ID:} \href{https://openalex.org/W3162214085}{W3162214085}  
   Esamina le differenze di genere nella gestione dell'immagine corporea sui social media, mostrando che i ragazzi tendono ad avere strategie di coping più attive rispetto alle ragazze (Mahon \& Hevey, 2021).

3. **Mitigating Harms of Social Media for Adolescent Body Image and Eating Disorders: A Review**  
   \textit{ID:} \href{https://openalex.org/W4400330793}{W4400330793}  
   Analizza strategie per mitigare gli effetti negativi dei social media sull'immagine corporea e i comportamenti alimentari, sottolineando la necessità di ulteriori ricerche per sviluppare approcci efficaci (Mazzeo et al., 2024).

4. **Self-esteem and social media addiction level in adolescents: The mediating role of body image**  
   \textit{ID:} \href{https://openalex.org/W4382312891}{W4382312891}  
   Mostra un legame negativo tra autostima e dipendenza dai social media, con l'immagine corporea che svolge un ruolo mediatorio in questa relazione (Colak et al., 2023).

5. **The Perfect Storm: A Developmental–Sociocultural Framework for the Role of Social Media in Adolescent Girls’ Body Image Concerns and Mental Health**  
   \textit{ID:} \href{https://openalex.org/W4285588124}{W4285588124}  
   Discute il contesto socioculturale che influisce sulle preoccupazioni legate all'immagine corporea tra le ragazze adolescenti in relazione ai social media (Choukas-Bradley et al., 2022).

\section{Interpretazione Prudente}
I risultati indicano che i contenuti dei social media possono significativamente influenzare la percezione dell'immagine corporea negli adolescenti, in particolare attraverso l'esposizione a ideali di bellezza irraggiungibili. Le differenze di genere influenzano anche le modalità di gestione di queste influenze, con i ragazzi che tendono a sviluppare strategie di coping più proattive. Inoltre, l'autostima gioca un ruolo critico, fungendo da mediatore tra dipendenza dai social media e percezione dell'immagine corporea.

\section{Limiti}
La complessità delle relazioni tra esposizione ai social media e percezione dell'immagine corporea rende difficile stabilire causalità dirette. Le differenze culturali e generazionali possono influenzare i risultati, limitando l'applicabilità delle conclusioni.

\section{Conclusione Provvisoria}
L'esposizione ai contenuti dei social media è associata a una percezione negativa dell'immagine corporea negli adolescenti. Tuttavia, ulteriori ricerche sono necessarie per chiarire i meccanismi e sviluppare interventi efficaci che possano mitigare i rischi associati.

\end{document}